% Part: first-order-logic
% Chapter: proof-systems
% Section: tableaux

\documentclass[../../../include/open-logic-section]{subfiles}

\begin{document}

\iftag{FOL}
      {\olfileid[pt]{fol}{prf}{tab}}
      {\olfileid[pt]{pl}{prf}{tab}}

\olsection{\usetoken{P}{tableau}}

Enquanto muitos sistemas de !!{derivation} operam com disposições de
!!{sentence}s, os !!{tableau}s operam com !!{signed formula}s.
!!^a{signed formula} é um par consistindo em um sinal de valor de
verdade ($\True$ ou $\False$) e !!a{sentence}
\[
\sFmla{\True}{!A} \text{ ou } \sFmla{\False}{!A}.
\]
!!^a{tableau} consiste em !!{signed formula}s dispostas em uma
árvore ramificada para baixo. Começa com várias \emph{hipóteses}
e continua com !!{signed formula}s que resultam de uma das
!!{signed formula}s acima dela pela aplicação de uma das regras de
inferência. Cada regra nos permite adicionar uma ou mais !!{signed
  formula}s ao final de um ramo, ou duas !!{signed formula}s lado a
lado---neste caso um ramo se divide em dois, com as duas !!{signed
  formula}s adicionadas formando os finais dos dois ramos.

Uma regra aplicada a uma !!{signed formula} complexa resulta na adição
de !!{signed formula}s que são sub-!!{formula}s imediatas. Elas vêm
em pares, uma regra para cada um dos dois sinais.  Por exemplo, a
regra $\TRule{\True}{\land}$ se aplica a $\sFmla{\True}{!A \land !B}$,
e permite a adição tanto das duas !!{signed formula}s
$\sFmla{\True}{!A}$ quanto de~$\sFmla{\True}{!B}$ ao final de qualquer ramo
que contém $\sFmla{\True}{!A \land !B}$, e a regra
$\TRule{\False}{!A \land !B}$ permite que um ramo seja dividido
adicionando $\sFmla{\False}{!A}$ e $\sFmla{\False}{!B}$ lado a lado.
!!^a{tableau} é fechado se cada um de seus ramos contém um
par correspondente de !!{signed formula}s $\sFmla{\True}{!A}$ e
$\sFmla{\False}{!A}$.

A relação $\Proves$ baseada em !!{tableau}s é definida da seguinte
forma: $\Gamma \Proves !A$ se, e somente se, há algum conjunto
finito~$\Gamma_0 = \{!B_1, \dots, !B_n\} \subseteq \Gamma$ tal que há
um !!{tableau} fechado para as hipóteses
\[
\{\sFmla{\False}{!A}, \sFmla{\True}{!B_1}, \dots, \sFmla{\True}{!B_n}\} 
\]
Por exemplo, aqui está um !!{tableau} fechado que mostra que $\Proves
(!A \land !B) \lif !A$:
\begin{oltableau}
  [\sFmla{\False}{(\formula{A} \land \formula{B}) \lif \formula{A}}, just = \TAss
    [\sFmla{\True}{\formula{A} \land \formula{B}}, just = {\TRule{\False}{\lif}[1]}
      [\sFmla{\False}{\formula{A}}, just = {\TRule{\False}{\lif}[1]}
        [\sFmla{\True}{\formula{A}}, just = {\TRule{\True}{\lif}[2]}
          [\sFmla{\True}{\formula{B}}, just = {\TRule{\True}{\lif}[2]}, close]
        ]
      ]
    ]
  ]
\end{oltableau}

Um conjunto $\Gamma$ é inconsistente no cálculo de !!{tableau} se há
um !!{tableau} fechado para hipóteses
\[
\{\sFmla{\True}{!B_1}, \dots, \sFmla{\True}{!B_n}\} 
\]
para algum $!B_i \in \Gamma$.

Os !!^{tableau}s foram inventados na década de 1950 independentemente
por Evert Beth e Jaakko Hintikka, e simplificados e popularizados por
Raymond Smullyan. São muito fáceis de usar, pois construir um
!!{tableau} é um procedimento muito sistemático. Por causa da natureza
sistemática dos !!{tableau}s, eles também se prestam à implementação
por computador. No entanto, um !!{tableau} é frequentemente difícil de
ler e sua conexão com provas às vezes não é fácil de ver. A abordagem
também é bastante geral, e muitas lógicas diferentes têm sistemas de
!!{tableau}. Os !!^{tableau}s também nos ajudam a encontrar
!!{structure}s que satisfazem !!{sentence}s (ou conjuntos de
!!{sentence}s) dadas: se o conjunto é satisfatível, ele não terá um
!!{tableau} fechado, isto é, qualquer !!{tableau} terá um ramo aberto.
A !!{structure} satisfatória pode ser ``lida'' de um ramo aberto,
desde que toda regra que seja possível aplicar tenha sido aplicada
naquele ramo.  Há também uma conexão muito estreita com o cálculo de
sequentes: essencialmente, um !!{tableau} fechado é uma
!!{derivation} condensada no cálculo de sequentes, escrita de cabeça
para baixo.

\end{document}
